\chapter{Pneumocystis Jirovecii PCR Test}

\section*{What Is This Test?}
Pneumocystis jirovecii PCR detects fungal DNA in a respiratory specimen and supports evaluation of Pneumocystis pneumonia (PJP), particularly in people with HIV, hematologic malignancy, organ transplantation, or immunosuppressive treatment.

\section*{Alternative Names}
\begin{itemize}
  \item PJP PCR
  \item Pneumocystis jirovecii PCR
  \item Pneumocystis DNA Test
  \item Pneumocystis jirovecii NAAT
\end{itemize}
\section*{How It Is Performed}
PCR is performed on an induced-sputum or lower-respiratory specimen. Quantitative assays may report a fungal DNA burden, but thresholds vary by assay and cannot be interpreted without the clinical context. Serum beta-D-glucan is often used as an adjunct, not as a replacement for respiratory testing.

\section*{Why It Is Ordered}
\begin{itemize}
  \item Evaluation of suspected PJP in an immunocompromised patient with respiratory symptoms or hypoxemia
  \item Supporting diagnosis when routine bacterial and viral testing does not explain the illness
  \item Investigation of breakthrough or recurrent pneumonia in a patient at risk for PJP
\end{itemize}

\section*{Interpretation and Limitations}
PCR is highly sensitive but may detect colonization rather than active pneumonia, especially in upper-airway specimens. A positive result must be interpreted with symptoms, oxygenation, imaging, microscopic staining or immunofluorescence when available, and beta-D-glucan. A negative result lowers the likelihood of PJP but does not exclude disease if the specimen is poor or fungal burden is low.
\section*{Specimen and Preparation}
Testing uses an induced-sputum or lower-respiratory specimen selected by the treating team; follow collection, transport, and infection-control instructions. Fasting is not required. Tell the clinician about HIV status, transplant or cancer treatment, corticosteroids, prophylaxis, oxygen requirement, and recent antimicrobial therapy.

\section*{Risks and Safety}
The PCR assay has no physical risk. Respiratory specimen collection may cause coughing or short-lived discomfort and is managed by the clinical team. Breathlessness, blue lips, confusion, or worsening oxygen levels require urgent care and should not await PCR results.

\section*{Understanding Results and Follow-up}
The report should state specimen type, qualitative result, and quantitative value or assay cutoff when available. PCR can detect colonization as well as disease, particularly in upper-airway samples, so a positive result is not sufficient by itself. Follow-up combines symptoms, oxygenation, respiratory findings, microscopy or staining when available, beta-D-glucan, and other diagnoses; a negative result may need reconsideration if the specimen is inadequate.

\section*{Regulatory Status and Authorization Date}
This chapter describes a test category, not a specific commercial product. FDA, CDSCO, EU IVDR, and MHRA approval or registration dates therefore cannot be assigned without the assay manufacturer, product name, instrument, and intended use. For laboratory-developed tests, an individual agency approval date may not exist. See the Preface for the interpretation of regulatory status.

\clearpage
